\section{이론}

실험을 정당화하는 데 긴 형식 체계가 필요한 것은 아니다. 필요한 것은 세 가지다. 첫째와 둘째는 구조적 명제다. 왜 stationary key-side 증폭이 순차적 coverage와 어긋나는지, 왜 query-side suppression이 올바른 history-free 근사인지를 설명한다. 마지막 정리는 perturbation 언어를 엄밀한 부등식으로 고정한다.

\paragraph{범위.}
이 절의 역할은 의도적으로 좁다. 두 개의 구조적 사실과 하나의 진단용 perturbation bound를 증명한다. 반대로, benchmark F1을 직접 증명하지 않으며, exact facet-history tracking도 주장하지 않는다.

\begin{table}[t]
\centering
\caption{본문 핵심 주장들의 증명 상태.}
\label{tab:proof-status-ko}
\small
\begin{tabularx}{\linewidth}{@{}Xlll@{}}
\toprule
주장 & 유형 & 완전 증명 여부 & 추가 모델링 가정 의존성 \\
\midrule
Stationary key-side steering은 explicit coverage variable을 도입하지 않는다 & 구조적 & 예 & 거의 없음 \\
음의 query-side projection은 coverage에 맞는 부호 의미를 가진다 & 구조적 & 예 & 거의 없음 \\
Attention output perturbation은 qaMSE와 $\rho^4$ remainder로 제어된다 & 진단용 & 예 & boundedness 가정 필요 \\
Q-bias가 MetaTool Subtask4의 second-tool recovery를 높인다 & 경험적 & 아니오 & 실험으로 보여야 함 \\
\bottomrule
\end{tabularx}
\end{table}

\paragraph{정리 1의 상시 가정.}
Perturbation 정리에는 다음의 분석 가정만 쓴다:
\begin{enumerate}
\item bounded query norm $\|q\| \le Q_{\max}$;
\item bounded value norm $\|v_t\| \le V_{\max}$;
\item row-wise key perturbation bound $\|e_t\| \le \rho$.
\end{enumerate}
앞의 두 명제는 이런 boundedness 가정에 의존하지 않는다.

\begin{proposition}{Stationary key-side steering은 explicit coverage variable을 도입하지 않는다}{}{}
선형사상 $P$와 스칼라 $\alpha$를 고정하고, 모든 디코딩 단계에서 같은 key-side 연산자
 \begin{equation*}
 k'=(I+\alpha P)k
 \end{equation*}
를 적용하자. 그러면 이 개입 자체는 어떤 facet이 이미 방출되었는지에 대한 정보를 담을 수 없다. 특히 현재 key tensor가 같다면, 과거에 어떤 facet이 사용되었는지와 무관하게 변환된 key는 완전히 동일하다.
\end{proposition}

\paragraph{해석.}
이 명제는 통계적이기보다 구조적이다. 시간에 대해 불변인 key 연산자는 single-target 벤치마크에서는 도움을 줄 수 있어도, ``이 facet은 이미 썼다''라는 explicit state를 스스로 추가하지는 못한다. 그런 정보는 key-side 연산자 바깥의 모델 상태에 이미 들어 있어야 한다. 부록 A에 짧은 증명을 실었다.

\paragraph{증명하지 않는 것.}
이 명제는 key-side steering이 항상 무의미하다고 말하지 않는다. 다만 stationary key-side map만으로는 순차적 멀티-툴 복원에 필요한 explicit coverage variable을 새로 도입하지 못한다고 말한다.

\begin{proposition}{Query-side suppression은 가장 단순한 history-free coverage 근사다}{}{}
온톨로지 부분공간으로의 직교사영을 $P_{\mathrm{ont}}$라고 하고
 \begin{equation*}
 q'=(I+\beta P_{\mathrm{ont}})q
 \end{equation*}
를 $\beta<0$에서 정의하자. 그러면 $q$의 ontology-aligned coefficient는 모두 균일하게 축소되고, ontology subspace에 직교한 성분은 그대로 유지된다. 따라서 음의 query-side bias는 외부 상태를 추가하지 않으면서도 지배적인 ontology-aligned mass의 반복 사용을 줄이는 가장 단순한 history-free 연산자다.
\end{proposition}

\paragraph{해석.}
이 명제가 exact facet tracking을 증명하는 것은 아니다. 더 좁고 충분한 사실만을 보여준다. history-free 선형 편집 가운데, 음의 query-side bias는 반복 facet 사용을 줄이는 데 맞는 부호 의미를 가진다. 실제 대수 전개는 부록 A에 정리했다.

\paragraph{증명하지 않는 것.}
이 명제는 모델이 항상 올바른 두 번째 도구로 전환된다는 것을 보장하지 않는다. 다만 외부 상태 없이 쓸 수 있는 history-free 선형 편집 가운데, 반복 facet 사용을 억제하는 부호 의미를 가진 가장 단순한 연산자를 식별한다.

\begin{theorem}{Attention-weighted perturbation bound}{}{}
질의 $q$, 원래 key 행렬 $K$, row-wise 오차 $\|e_t\|\le\rho$를 갖는 교란 key 행렬 $\hat K=K+E$, 그리고 대응하는 attention output $\hat o(q), o(q)$를 두자. 또한
 \begin{equation*}
 \begin{aligned}
 \alpha_t(q) &= \frac{q\cdot e_t}{\sqrt d},\\
 \mathrm{qaMSE}(q;E) &= \sum_{t=1}^{T}s_t(q)\,\alpha_t(q)^2
 \end{aligned}
 \end{equation*}
를 정의하자. 여기서 $s_t(q)$는 원래 attention weight다. bounded-query, bounded-value 가정 아래에서
 \begin{equation*}
 C_1=\frac{2Q_{\max}^4V_{\max}^2}{d^2}
 \end{equation*}
라는 명시적 상수가 존재하여
 \begin{equation*}
 \mathbb{E}_{q}\|\hat o(q)-o(q)\|^2
 \le
 2\,\mathbb{E}_{q}\!\left[\mathrm{qaMSE}(q;E)\,\mathrm{Var}_{s}[V](q)\right]
 + C_1\rho^4
 \end{equation*}
가 성립한다.
\end{theorem}

\paragraph{해석.}
엄밀한 부분은 이 정리 자체이고, 경험적 검증은 그 다음 단계다. 이 정리는 두 항으로 읽으면 이해가 쉽다. 첫 번째 항은 attention-weighted 1차 기여이고, 두 번째 항은 $\rho^4$로 제어되는 고차 remainder다. 표~\ref{tab:bound-ko}는 측정된 perturbation-to-bound ratio의 중앙값이 Qwen에서는 $2.36\times 10^{-8}$, Llama에서는 $6.37\times 10^{-8}$임을 보여준다. 이 정리가 end-to-end accuracy를 증명하는 것은 아니지만, Q-side와 K-side 개입을 비교할 때 perturbation 언어를 보수적인 진단 도구로 쓰는 것은 정당화한다. 증명은 부록 A에 네 단계로 나누어 제시한다.

\paragraph{증명하지 않는 것.}
이 정리는 benchmark theorem이 아니다. qaMSE가 더 작으면 반드시 tool-selection F1이 더 좋다고 말하는 것이 아니라, 주어진 boundedness 가정 아래에서 key 오차가 attention output을 얼마나 교란하는지를 제어하는 진단식을 준다.

\begin{remark}{}{}{}
평가한 연산자는 history-free다. ontology subspace로의 사영을 억제할 뿐, 이미 방출된 facet을 명시적으로 추적하지는 않는다. 따라서 이론의 역할도 의도적으로 좁다. 정확한 step-adaptive control을 형식화하는 것이 아니라, 개입의 부호, 특이성, perturbation 안정성을 해석하는 데 초점을 둔다.
\end{remark}
